%%
%% Copyright 2007-2020 Elsevier Ltd
%%
%% This file is part of the 'Elsarticle Bundle'.
%% ---------------------------------------------
%%
%% It may be distributed under the conditions of the LaTeX Project Public
%% License, either version 1.2 of this license or (at your option) any
%% later version.  The latest version of this license is in
%%    http://www.latex-project.org/lppl.txt
%% and version 1.2 or later is part of all distributions of LaTeX
%% version 1999/12/01 or later.
%%
%% The list of all files belonging to the 'Elsarticle Bundle' is
%% given in the file `manifest.txt'.
%%

%% Template article for Elsevier's document class `elsarticle'
%% with numbered style bibliographic references
%% SP 2008/03/01
%%
%%
%%
%% $Id: elsarticle-template-num.tex 190 2020-11-23 11:12:32Z rishi $
%%
%%
\documentclass[preprint,12pt, compress]{elsarticle}

%% Use the option review to obtain double line spacing
%%\documentclass[authoryear,preprint,review,12pt]{elsarticle}

%% Use the options 1p,twocolumn; 3p; 3p,twocolumn; 5p; or 5p,twocolumn
%% for a journal layout:
%% \documentclass[final,1p,times]{elsarticle}
%% \documentclass[final,1p,times,twocolumn]{elsarticle}
%% \documentclass[final,3p,times]{elsarticle}
%%b\documentclass[final,3p,times,twocolumn]{elsarticle}
%% \documentclass[final,5p,times]{elsarticle}
%% \documentclass[final,5p,times,twocolumn]{elsarticle}

%% For including figures, graphicx.sty has been loaded in
%% elsarticle.cls. If you prefer to use the old commands
%% please give \usepackage{epsfig}

%% The amssymb package provides various useful mathematical symbols
\usepackage{amssymb}
\usepackage{hyperref}
%% The amsthm package provides extended theorem environments
%% \usepackage{amsthm}

%% Packages added by Lenz
\usepackage{xcolor}
\usepackage{listings}
\usepackage{bm}
\usepackage{relsize}
\usepackage{tabularx}
\usepackage{booktabs}

\usepackage{amsmath}
\usepackage{braket}

\newcommand{\TODO}[1]{\textcolor{red}{#1}}
\newcommand{\ionion}{_\mathrm{ii}}
\newcommand{\eleion}{_\mathrm{ei}}
\newcommand{\eleele}{_\mathrm{ee}}
\newcommand{\ion}{_\mathrm{i}}
\newcommand{\ele}{_\mathrm{e}}
\newcommand{\br}{\bm{r}}
\newcommand{\bR}{\bm{R}}
\newcommand{\ubr}{\underline{\bm{r}}}
\newcommand{\ubR}{\underline{\bm{R}}}
\newcommand{\bu}{{\bf u}}
\newcommand{\bU}{{\bf U}}
\newcommand{\bF}{{\bf F}}
\newcommand{\bmmatrix}[1]{\bm{\mathrm{#1}}}
\newcommand{\bk}{\bm{k}}
\newcommand{\s}{_\mathrm{{\scriptscriptstyle S}}}
\newcommand{\h}{_\mathrm{{\scriptscriptstyle H}}}
\newcommand{\xc}{_\mathrm{{\scriptscriptstyle XC}}}
\newcommand{\hxc}{_\mathrm{{\scriptscriptstyle HXC}}}
\newcommand{\Hartree}{_\mathrm{{\scriptscriptstyle H}}}
\newcommand{\kB}{k_\mathrm{B}}

\DeclareMathOperator{\Tr}{Tr}


\journal{Computer Physics Communications}
\newcolumntype{b}{X}
\newcolumntype{s}{>{\hsize=.4\hsize}X}

\begin{document}

\begin{frontmatter}

%% Title, authors and addresses

%% use the tnoteref command within \title for footnotes;
%% use the tnotetext command for theassociated footnote;
%% use the fnref command within \author or \address for footnotes;
%% use the fntext command for theassociated footnote;
%% use the corref command within \author for corresponding author footnotes;
%% use the cortext command for theassociated footnote;
%% use the ead command for the email address,
%% and the form \ead[url] for the home page:
%% \title{Title\tnoteref{label1}}
%% \tnotetext[label1]{}
%% \author{Name\corref{cor1}\fnref{label2}}
%% \ead{email address}
%% \ead[url]{home page}
%% \fntext[label2]{}
%% \cortext[cor1]{}
%% \affiliation{organization={},
%%             addressline={},
%%             city={},
%%             postcode={},
%%             state={},
%%             country={}}
%% \fntext[label3]{}

\title{Working Title: MANDALA Conceptual Paper}

%% use optional labels to link authors explicitly to addresses:
%% \author[label1,label2]{}
%% \affiliation[label1]{organization={},
%%             addressline={},
%%             city={},
%%             postcode={},
%%             state={},
%%             country={}}
%%
%% \affiliation[label2]{organization={},
%%             addressline={},
%%             city={},
%%             postcode={},
%%             state={},
%%             country={}}

\author[hzdr,casus]{Bartosz Brzoza}

\author[hzdr,casus]{Attila Cangi}%\corref{cor1}}


\affiliation[hzdr]{organization={Helmholtz-Zentrum Dresden-Rossendorf},%Department and Organization
            %addressline={},
            city={Dresden},
            postcode={01328},
            %state={},
            country={Germany}}

\affiliation[casus]{organization={Center for Advanced Systems Understanding},%Department and Organization
            %addressline={},
            city={Görlitz},
            postcode={02826},
            %state={},
            country={Germany}}

\affiliation[tud]{organization={Institute of Artificial Intelligence, Dresden University of Technology},%Department and Organization
            %addressline={},
            city={Dresden},
            postcode={01069},
            %state={},
            country={Germany}}

%\cortext[cor1]{Corresponding author: a.cangi@hzdr.de}


\begin{abstract}
We present ...
\end{abstract}


%%Graphical abstract
%\begin{graphicalabstract}
%\includegraphics{grabs}
%\end{graphicalabstract}

%%Research highlights
%\begin{highlights}
%\item Research highlight 1
%\item Research highlight 2
%\end{highlights}

\begin{keyword}
Machine Learning  \sep Neural Networks \sep Electronic Structure Theory \sep Density Functional Theory
%% keywords here, in the form: keyword \sep keyword

%% PACS codes here, in the form: \PACS code \sep code

%% MSC codes here, in the form: \MSC code \sep code
%% or \MSC[2008] code \sep code (2000 is the default)

\end{keyword}

\end{frontmatter}

{\bf PROGRAM SUMMARY/NEW VERSION PROGRAM SUMMARY}
  %Delete as appropriate.

% LF: This is a required section/field for CPiC papers
\begin{small}
\noindent
{\em Program Title: }MANDALA (\texttt{MANDALA})       \\
{\em CPC Library link to program files:} (to be added by Technical Editor) \\
{\em Developer's repository link:} \url{https://github.com/XXX} \\
{\em Code Ocean capsule:} (to be added by Technical Editor)\\
{\em Licensing provisions(please choose one):} XXX (BSD 3-clause)  \\
{\em Programming language:} Python                                   \\
{\em Supplementary material:} None                                 \\
{\em Nature of problem(approx. 50-250 words):} XXX\\
{\em Solution method(approx. 50-250 words):} XXX\\
{\em Additional comments including restrictions and unusual features (approx. 50-250 words):} None\\
   \\

\begin{thebibliography}{0}
\bibitem{1} MANDALA (Version 1.X.X) [Computer software]. https://doi.org/XX.XXXX/XXXXX

% Program Summary section.

\end{thebibliography}
\end{small}


%% main text
\section{Introduction}
\label{s.introduction}

XXX


\section{Theoretical background}
\label{s.background}

\subsection{Electronic structure problem and density functional theory}
XXX

\subsection{Learning the electronic structure problem with E3-equivariant graph neural netoworks}
XXX




\section{Numerical implementation}
\label{s.implementation}

\subsection{Code structure}
\label{ss.code.structure}


\subsection{Data sampling}
\label{ss.data.sampling}



\subsection{Hyperparameter optimization and model training}
\label{ss.hyperparameters}


\subsection{Observable calculation}
\label{ss.observables}


\subsection{Parallelization and hardware acceleration}
\label{ss.parallelization}



\subsection{Data management}
\label{ss.data.management}





\section{Numerical benchmarks}
\label{s.benchmarks}






\section{Conclusion}
\label{s.conclusion}




\section*{Acknowledgments}

This work was partially supported by the Center for Advanced Systems Understanding (CASUS) which is financed by Germany's Federal Ministry of Education and Research (BMBF) and by the Saxon state government out of the State budget approved by the Saxon State Parliament. Computations were performed on a Bull Cluster at the Center for Information Services and High-Performance Computing (ZIH) at Technische Universit\"at Dresden and on the cluster Hemera at Helmholtz-Zentrum Dresden-Rossendorf (HZDR).

%% The Appendices part is started with the command \appendix;
%% appendix sections are then done as normal sections
\appendix

\section{Computational workflow example}
\label{s.workflow}




%% Statement on the use of generative AI
\vspace{1cm}
XXX


%% If you have bibdatabase file and want bibtex to generate the
%% bibitems, please use
%%
\bibliographystyle{elsarticle-num}
\bibliography{references.bib}

%% else use the following coding to input the bibitems directly in the
%% TeX file.

%% \begin{thebibliography}{00}

%% \bibitem{label}
%% Text of bibliographic item

%% \bibitem{}

%% \end{thebibliography}
\end{document}
%% \endinput
%%
%% End of file `elsarticle-template-num.tex'.
